% ===========================================================================
%  sviluppatore/02-ambiente/03-dati.tex
% ===========================================================================
\section{Dove stanno i dati}\label{sez:dati}

Tutti i percorsi del sistema sono definiti in un file solo,
\file{core/paths.py}, e nessun altro modulo costruisce percorsi a mano. Chi
deve aggiungere una cartella la aggiunge lì.

\subsection{La cartella di lavoro}

La radice non è necessariamente quella del repository. \simbolo{PROJECT\_ROOT}
si decide a ogni avvio, in questo ordine:

\begin{enumerate}
  \item \cod{LYRA\_HOME} impostata nell'ambiente \freccia{} si usa quella. È il
        modo per tenere i dati su un disco diverso, o condividerli fra più
        installazioni.
  \item Applicazione impacchettata (PyInstaller, Electron) \freccia{} la
        cartella utente del sistema operativo:
        \file{\%LOCALAPPDATA\%/Lyra} su Windows,
        \file{\textasciitilde/Library/Application Support/Lyra} su macOS,
        \file{\textasciitilde/.local/share/Lyra} su Linux.
  \item Altrimenti \freccia{} la radice del repository.
\end{enumerate}

\begin{perche}{Perché i dati escono dal bundle}
In un'applicazione impacchettata il codice vive in una cartella di sola
lettura che ogni aggiornamento sovrascrive. Se i \file{parquet} e i risultati
dei backtest stessero lì dentro, un aggiornamento cancellerebbe mesi di
lavoro. \simbolo{CODE\_ROOT} --- dove sta il codice --- e
\simbolo{PROJECT\_ROOT} --- dove si scrive --- sono due cose distinte proprio
per questo, e vanno tenute distinte anche quando coincidono, come in sviluppo.
\end{perche}

\subsection{Che cosa c'è sotto \file{data/}}

\begin{center}
\begin{tabular}{@{}L{40mm}L{104mm}@{}}
\toprule
\textbf{Cartella} & \textbf{Contenuto} \\
\midrule
\rigaTenue
\file{data/assets/<TF>/}
  & OHLCV, un \file{parquet} per simbolo, divisi per timeframe:
    \file{D1/}, \file{H4/}, \file{H1/}, \file{M30/}, \file{M15/},
    \file{M5/}. Lo stesso simbolo può esistere a più timeframe, ciascuno
    con il proprio intervallo di date. \\
\file{data/driver/}
  & Le serie che guidano i regimi (SPY, DXY, VIX\ldots). Stanno a parte
    perché possono avere un timeframe diverso dagli asset. \\
\rigaTenue
\file{data/rates/}      & Tassi di politica monetaria per valuta, da FRED. \\
\file{data/news/}       & Archivio delle notizie. \\
\rigaTenue
\file{data/earnings/}   & Date e sorprese degli utili, per simbolo. \\
\file{data/fundamentals/} & Bilanci, per simbolo. \\
\rigaTenue
\file{data/mms/}        & Feature della heatmap di liquidità, per
                          (simbolo, timeframe). \\
\file{data/\_universe.json}
  & Che cosa il sistema sa di avere. Non è un indice generato: è il
    registro, e una cartella piena di \file{parquet} non censiti qui per il
    sistema è vuota. \\
\bottomrule
\end{tabular}
\end{center}

\noindent
I risultati stanno altrove, sotto \file{output/}: \file{output/strategies/}
per gli stream (è da qui che la dashboard li legge, seguendo il
\cod{output\_path\_pattern} dichiarato da ciascun artefatto),
\file{output/ml/} per feature e predizioni, \file{output/reports/} per i
report aggregati.

\begin{nota}{I percorsi legacy sono ancora letti}
\file{tests/outputs/}, \file{data/feature\_store/}, \file{data/ml\_outputs/}
sono posizioni storiche. L'helper \simbolo{discover\_parquet()} cerca prima
nei percorsi nuovi e poi in questi, così una installazione vecchia continua a
funzionare. Scrivere lì non è più previsto.
\end{nota}

\subsection{Portare dentro i primi dati}

Gli script di scarico stanno in \file{scripts/} e sono uno per sorgente:

\begin{codiceInterno}{console}{I quattro che servono per cominciare}{scripts/}
python scripts/download_d1.py              # azioni ed ETF, giornaliero
python scripts/download_regime_drivers.py  # le serie dichiarate dai mercati di regime
python scripts/download_fred_rates.py      # tassi (serve FRED_API_KEY)
python scripts/download_polygon_m5.py      # intraday (serve POLYGON_API_KEY)
\end{codiceInterno}

\noindent
Lo stesso si fa dall'interfaccia, in \voce{Data}, che chiama le rotte
\cod{/api/data/download/<modalità>} e mostra l'avanzamento. La differenza è
che dall'interfaccia il lavoro parte come \term{job} con un registro
consultabile, dalla riga di comando no.

\begin{attenzione}{\cod{download\_regime\_drivers.py} legge i file di mercato}
Non ha un elenco di simboli suo: scarica i \cod{driver} dichiarati nei file
sotto \file{regime/markets/}. Se si aggiunge un mercato che guarda una serie
nuova, questo script è il modo di procurarsela --- e finché non lo si lancia,
quel mercato non modula niente e lo dice solo in un avviso nel registro.
\end{attenzione}
